%%%%%%%%%%%%%%%%%%%%%%%%%%%%%%%%%%%%%%%%%%%%%%%%%%%%%%%%%%%%%%%%%%%%%%%%
% Plantilla TFG/TFM
% Universidad de Murcia. Facultad de Informática
%%%%%%%%%%%%%%%%%%%%%%%%%%%%%%%%%%%%%%%%%%%%%%%%%%%%%%%%%%%%%%%%%%%%%%%%


\chapter*{Extended Abstract}

This Final Degree Project presents the comprehensive design, implementation, and rigorous evaluation of an intelligent conversational agent specialized in the cinematographic domain. The primary objective is to bridge the significant technical gap between the natural language comprehension capabilities of modern Large Language Models (LLMs) and the factual, deterministic precision inherent to Semantic Web Knowledge Graphs. By synergizing these two distinct technological paradigms, the project directly addresses one of the most critical and widely discussed challenges in modern generative Artificial Intelligence: the mitigation and elimination of "hallucinations"—the generation of plausible yet entirely fabricated information.

\section*{1. Introduction and Motivation}
In recent years, LLMs have revolutionized the way users and developers interact with information systems, enabling highly fluent conversational interfaces. However, their stochastic nature makes them inherently unreliable for querying specific, factual, or specialized data, as they are prone to generating incorrect answers with high confidence. On the other hand, Knowledge Graphs and the Semantic Web provide impeccably structured, verifiable data environments accessible via precise query languages like SPARQL. The steep learning curve and strict syntax of SPARQL prevent non-expert users from accessing this immense wealth of information. This project proposes a robust hybrid solution: an autonomous conversational agent that acts as a transparent natural language interface for a local Knowledge Graph, intelligently translating user queries into SPARQL, executing them locally, and returning the exact data in a human-readable format.

\section*{2. State of the Art}
The theoretical foundation of this work rests firmly on three technological pillars: the Semantic Web, Large Language Models, and Retrieval-Augmented Generation (RAG) paradigms. The Semantic Web, formalized and standardized by the W3C, relies heavily on the Resource Description Framework (RDF) to elegantly model complex information as subject-predicate-object triples. Wikidata, a massive, multilingual, and collaborative knowledge base, serves as the primary and authoritative data source for this project.

To effectively overcome the static knowledge limitations and hallucination issues of standard LLMs, RAG architectures dynamically retrieve external data and inject it into the model's contextual window. This project pushes the traditional RAG paradigm significantly further by implementing an "Agentic AI" workflow using the LangChain framework. Within this architecture, the LLM is not relegated to merely retrieving text; it actively reasons, plans, writes code, tests queries against the database, and refines its approach.

\section*{3. System Architecture and Methodology}
Given the highly experimental and empirical nature of advanced prompt engineering, an agile, iterative, and test-driven methodology was adopted throughout the development lifecycle. The overall system architecture is modularly divided into three highly cohesive main layers:
\begin{itemize}
    \item \textbf{Presentation Layer:} A highly interactive and responsive graphical user interface built entirely with Streamlit. It provides a seamless chat-like experience while simultaneously displaying the internal reasoning process (Chain-of-Thought) and the generated SPARQL code to the user, ensuring maximum transparency.
    \item \textbf{Logic and Agent Layer:} The cognitive core of the application, powered by LangChain and Llama3 running locally via Ollama. This layer handles the complex tasks of natural language understanding, entity extraction, Wikidata term disambiguation, and the core iterative SPARQL generation loop.
    \item \textbf{Data Layer:} A highly optimized, in-memory RDF local repository serialized in \texttt{.ttl} format, managed by the RDFLib Python library. It contains a densely populated subset of domain-specific cinematic data.
\end{itemize}

\section*{4. Local Knowledge Graph Construction}
Relying directly on public SPARQL endpoints (such as the Wikidata Query Service) during real-time user interaction introduces unacceptable latency and severe risks of being blocked by strict API rate limits. Therefore, a custom, highly tailored data extraction pipeline was developed. The system programmatically queries Wikidata entirely offline to extract a densely connected, self-contained subgraph focusing specifically on movies, directors, actors, genres, and release dates. This data is then serialized into the \textit{Turtle} format. This approach provides a lightning-fast, secure, and fully controlled environment for the agent to execute queries without relying on unstable external network dependencies.

\section*{5. The Reflexion Loop}
Arguably the most innovative and impactful aspect of the system's implementation is the agent's autonomous, self-healing error-correction mechanism. Translating fluid natural language into strict, unforgiving SPARQL syntax is a heavily error-prone task for any LLM. To guide the model, the agent is provided with a strict, curated dictionary of allowed RDF properties (e.g., \texttt{wdt:P57} exclusively for "director"). 

When the agent confidently generates a SPARQL query, the Logic Layer attempts to execute it locally against the Data Layer via RDFLib. If a syntax or execution error occurs, the process does not fail catastrophically. Instead, the runtime traceback is intercepted, parsed, and fed back to the LLM in a specialized "reflexion prompt". This prompt instructs the model to critically analyze its own previous mistake, understand the syntax error, and autonomously rewrite the query. This iterative self-correction loop drastically increases the overall robustness, reliability, and success rate of the system.

\section*{6. Evaluation and Results}
To quantitatively validate the system's performance and viability, a comprehensive automated evaluation framework was built. A rigorous dataset of 61 complex natural language queries was carefully created, covering a wide array of semantic relationships and designed specifically to compare the performance of various large language models. The chosen baseline model for this project (Llama3) achieved an outstanding 90.16\% accuracy rate in the crucial steps of entity extraction and disambiguation. Furthermore, it also achieved a 90.16\% overall success rate in generating valid, executable SPARQL queries that returned the factually correct answer.

Within this framework, we also rigorously evaluated the execution latency per question. Llama3 demonstrated its ideal suitability for real-time conversational AI applications by achieving a remarkable average execution time of just 3.66 seconds per question, which was comfortably the lowest latency among all the sophisticated models tested.

An in-depth analysis of the remaining failure cases revealed that most issues stemmed from extreme idiomatic ambiguities in the user's phrasing or minor, residual API rate limits encountered during the entity disambiguation phase, rather than fundamental logical reasoning failures by the LLM itself. The extraordinarily high success rate in autonomous query generation conclusively proves that heavily constraining the LLM's vast search space through strict prompt engineering and predefined property dictionaries is a highly effective, production-ready strategy.

\section*{7. Conclusions}
The empirical results confirm beyond doubt that integrating powerful LangChain agents with strongly-typed, local Knowledge Graphs is a highly viable and superior technical solution for building transparent, hallucination-free Question-Answering systems. The resulting agentic workflow successfully and completely hid the underlying complexity of the Semantic Web from the end-user, all while retaining and leveraging its absolute factual precision. Future work will strategically focus on improving the semantic similarity algorithms to ensure flawless entity disambiguation, and expanding the underlying ontology to comfortably support highly complex, multi-hop reasoning questions.

\vspace{1cm}
\textbf{Keywords:} Semantic Web, Knowledge Graphs, SPARQL, Large Language Models, Llama3, RAG, LangChain, Agentic AI, Hallucination Mitigation, Streamlit.
